
\section{Evaluation}




\subsection{Dataset Construction}

Dataset A.

Dataset B.



\subsection{Experiment Setup}

\begin{itemize}
    % baselines 要和哪些方法做对比
    \item \textbf{Accuracy Evaluation (RQ1)}: How is the overall accuracy of \tool compared with the state-of-the-arts? 
    \item \textbf{Effectiveness Evaluation (RQ2)}: How is the effectiveness of \tool in different API compatibility types? % 证明不同API 类型上，准确率都还行
    
    \item \textbf{LLM Generalizability (RQ3)}: How is the accuracy of \tool using different LLMs?
    \item \textbf{Ablation Study (RQ4)}: 
\end{itemize}

\textbf{Baselines. } We selected XXX~\cite{zhang2022fldetector}, XXX~\cite{zhang2022fldetector}, and XXX~\cite{zhang2022fldetector} as the comparison.

\textbf{Metrics.} We use precision and recall as metrics, which are defined as follows:



\subsection{Results}



\subsubsection{Accuracy Evaluation (RQ1)}

% 通过实验证明“ 4. 现有对于API兼容性问题的检测的工作，主要的方法，\cite{A,  B, C-shenkang} 以及存在的局限性”

\subsubsection{RQ2}

\subsubsection{RQ3}

\subsubsection{RQ4}



